% Part: first-order-logic
% Chapter: beyond
% Section: other-logics

\documentclass[../../../include/open-logic-section]{subfiles}

\begin{document}

\olfileid{fol}{byd}{oth}

\olsection{Andere logica's}

Zoals u inmiddels wellicht hebt opgemerkt, is het niet moeilijk een nieuwe
logica te ontwerpen. Ook u kunt een eigen syntaxis kiezen, een deductief systeem
opstellen en daar een semantiek bij ontwerpen. Er kan enige vindingrijkheid
nodig zijn om het !!{derivation} systeem volledig te maken ten opzichte van de
semantiek, en het kan enige moeite kosten de rest van de wereld ervan te
overtuigen dat uw logica werkelijk interessant is. Daar staat tegenover dat u
urenlang verantwoord plezier kunt beleven aan het onderzoek van de wiskundige
en computationele eigenschappen ervan.

De afgelopen decennia is het aantal formele logica's explosief gegroeid. Fuzzy
logic is bedoeld om redeneringen over vage eigenschappen te modelleren.
Probabilistische logica is bedoeld om redeneringen over onzekerheid te
modelleren. Defaultlogica's en niet-monotone logica's zijn bedoeld om
weerlegbare vormen van redeneren te modelleren: ``redelijke'' gevolgtrekkingen
die door nieuwe informatie later ongedaan kunnen worden gemaakt. Er zijn
epistemische logica's, bedoeld om redeneringen over kennis te modelleren;
causale logica's, bedoeld om redeneringen over causale verbanden te modelleren;
en zelfs ``deontische'' logica's, bedoeld om redeneringen over morele en
ethische verplichtingen te modelleren. Al naargelang de voornaamste reden om
zulke systemen in te voeren filosofisch, wiskundig of computationeel is, treft
men deze systemen aan onder de noemer wiskundige logica, filosofische logica,
kunstmatige intelligentie, cognitiewetenschap of elders.

De lijst gaat maar door en de mogelijkheden lijken eindeloos. Misschien zullen
we Leibniz' droom om de gehele menselijke rede tot berekening terug te brengen
nooit verwezenlijken---maar dat hoeft ons er niet van te weerhouden het te
proberen.

\end{document}
